\section{The results of \emph{Algorithmic Agnosticism}}\label{sec:aa}

This section restates the parts of \emph{Algorithmic Agnosticism} (AA) that this paper uses, in the language of
philosophy of religion. The proofs are there.

\medskip\noindent\textbf{AA-D. Classes.}\phantomsection\label{item:AAD}
\citep[§2]{taylor2026aa}
\begin{itemize}
  \item \emph{Evidence} is what a person has observed, experienced, remembered or written down.
        It is non-factive about the world: ``it seemed that God spoke'', not ``God spoke''.
  \item An \emph{account} is a statement of how the world is.
  \item A \emph{judgment} is a way of passing from evidence to an account.
  \item An \emph{imposed constraint} is a criterion for choosing among accounts that the
        evidence does not determine.
  \item \emph{Truths of reason} (this paper's term) are logic, mathematics, and conceptual
        analysis.
\end{itemize}

\medskip\noindent\textbf{AA-P1. Finite evidence.}\phantomsection\label{item:AAP1}
\citep[§2]{taylor2026aa} Every person's evidence is finite.

\medskip\noindent\textbf{AA-R1. No establishment.}\phantomsection\label{item:AAR1}
\citep[§3]{taylor2026aa}
\begin{itemize}
  \item Every finite body of evidence fits infinitely many accounts that disagree about what has
        not yet been observed.
  \item No judgment of any kind establishes one of them as the true account: any judgment that selects one as
        the true account is wrong in some account consistent with the evidence.
\end{itemize}

\medskip\noindent\textbf{AA-R2. Permanent underdetermination.}\phantomsection\label{item:AAR2}
\citep[§6]{taylor2026aa} Accounts that agree on all possible evidence are never separated by any
judgment, even by an ideal judge with all the evidence there could ever be.

\medskip\noindent\textbf{AA-R3. What remains is choice.}\phantomsection\label{item:AAR3}
\citep[§§4--5, 9]{taylor2026aa} Among the accounts that the evidence leaves open, selection is
made by an imposed constraint, and the evidence does not determine which.
